% from ega2-1-fr.tex
\oldpage[II]{184}
l'injection canonique, on pose
$\mathcal{L}=j^*(\mathcal{J})=
\mathcal{J}\otimes_{\mathcal{O}_V}\mathcal{O}_X$, de sorte que
$j_*(\mathcal{L})=\mathcal{J}/\mathcal{J}^2$. On suppose que le morphisme
structural $p:X\to Y$ est séparé et quasi-compact, et que les conditions
suivantes sont remplies~:
\begin{enumerate}
  \item[\rm{(i)}] Il existe un $Y$-morphisme de type fini $\pi:V\to X$ tel
  que $\pi\circ j=1_X$, et par suite
  $\pi_*(\mathcal{J}/\mathcal{J}^2)=\mathcal{L}$.
  \item[\rm{(ii)}] Il existe un homomorphisme de $\mathcal{O}_X$-Modules
  \[
    \varphi:\mathcal{L}\longrightarrow
    \varprojlim_n\pi_*(\mathcal{J}/\mathcal{J}^{n+1})
  \]
  tel que le composé
  \[
    \mathcal{L}\xrightarrow{\varphi}
    \varprojlim_n\pi_*(\mathcal{J}/\mathcal{J}^{n+1})
    \xrightarrow{\alpha}\pi_*(\mathcal{J}/\mathcal{J}^2)=\mathcal{L}
  \]
  (où $\alpha$ est l'homomorphisme canonique) soit l'identité.
  \item[\rm{(iii)}] Il existe un $Y$-préschéma $C$, une $Y$-section
  $\varepsilon$ de $C$ et un $Y$-morphisme $q:V\to C$ tels que le diagramme
  \begin{equation}
  \label{II.8.9.4.1-fr}
    \xymatrix{
      X\ar[r]^{j}\ar[d]_{p} & V\ar[d]^{q}\\
      Y\ar[r]_{\varepsilon} & C
    }
  \tag{8.9.4.1}
  \end{equation}
  soit commutatif.
  \item[\rm{(iv)}] La restriction de $q$ à $W=V-j(X)$ est une immersion
  ouverte quasi-compacte dans $C$, dont l'image ne rencontre pas
  $\varepsilon(Y)$.
\end{enumerate}
Dans ces conditions, $\mathcal{L}$ est ample relativement à $p$.
\end{proposition}

\subsection{Le critère d'amplitude de Grauert : démonstration}
\label{subsection:II.8.10-fr}

\begin{lemma}[8.10.1]
\label{II.8.10.1-fr}
Soient $\pi:V\to X$ un morphisme, $j:X\to V$ une $X$-section de $V$ qui soit
une immersion fermée, $\mathcal{J}$ l'Idéal quasi-cohérent de
$\mathcal{O}_V$ définissant le sous-préschéma fermé de $V$ associé à $j$.
\begin{enumerate}
  \item[\rm{(i)}] Pour tout $n\geqslant 0$,
  $\pi_*(\mathcal{O}_V/\mathcal{J}^{n+1})$ et
  $\pi_*(\mathcal{J}/\mathcal{J}^{n+1})$ sont des
  $\mathcal{O}_X$-Modules quasi-cohérents, et on a
  $\pi_*(\mathcal{O}_V/\mathcal{J})=\mathcal{O}_X$,
  $\pi_*(\mathcal{J}/\mathcal{J}^2)=j^*(\mathcal{J})$.
  \item[\rm{(ii)}] Si $X=\{\xi\}=\operatorname{Spec}(k)$, où $k$ est un
  corps, $\varprojlim_n\pi_*(\mathcal{O}_V/\mathcal{J}^{n+1})$ est isomorphe
  au séparé complété de l'anneau local $\mathcal{O}_{j(\xi)}$ pour la
  topologie $\mathfrak{m}_{j(\xi)}$-préadique.
  \item[\rm{(iii)}] Supposons que $\mathcal{J}$ soit un
  $\mathcal{O}_V$-Module inversible (ce qui entraîne que
  \[
    \mathcal{L}=j^*(\mathcal{J})=
    \pi_*(\mathcal{J}/\mathcal{J}^2)
  \]
  est un $\mathcal{O}_X$-Module inversible) et qu'il existe un homomorphisme
  \[
    \varphi:\mathcal{L}\longrightarrow
    \varprojlim_n\pi_*(\mathcal{J}/\mathcal{J}^{n+1})
  \]
  tel que le composé
  \[
    \mathcal{L}\xrightarrow{\varphi}
    \varprojlim_n\pi_*(\mathcal{J}/\mathcal{J}^{n+1})
    \xrightarrow{\alpha}\pi_*(\mathcal{J}/\mathcal{J}^2)
  \]
  ($\alpha$ homomorphisme canonique) soit l'identité. Si on pose
  $\mathcal{S}=\bigoplus_{n\geqslant0}\mathcal{L}^{\otimes n}$, on déduit
  alors canoniquement de $\varphi$ un isomorphisme de
  $\mathcal{O}_X$-Algèbres du complété $\widehat{\mathcal{S}}$ de
  $\mathcal{S}$ relatif à sa filtration canonique (complété qui est isomorphe
  au produit $\prod_{n\geqslant0}\mathcal{L}^{\otimes n}$) sur
  $\varprojlim_n\pi_*(\mathcal{O}_V/\mathcal{J}^{n+1})$.
\end{enumerate}
\end{lemma}

\begin{proof}
Notons en premier lieu que le support du $\mathcal{O}_V$-Module
$\mathcal{O}_V/\mathcal{J}^{n+1}$ est $j(X)$ et celui de
$\mathcal{J}/\mathcal{J}^{n+1}$ est contenu dans $j(X)$. Dans le cas (ii),
$j(X)$ est un point fermé $j(\xi)$ de $V$,

